% !PRINTED 422
\weberpage{422}
%<<CONT
von $F_n[j(\tau), j(\tau)]$ nach Potenzen von $q=e^{\pi i \tau}$ mit $Cq^{-2N}$
[§\,69, (4)].

Wenn wir also die Faktoren der rechten Seite von (12) in gleicher Weise
entwickeln, so können wir sowohl $C$ als $N$ bestimmen.

Für einen Faktor der rechten Seite von (1) % [sic] Der Druck nennt (1); der vorangehende Satz und der angegebene Faktor zeigen, daß (12) gemeint ist.
\[
j(\tau) - j\left(\frac{c+\partial\tau}{a}\right)
\]
haben wir folgende Anfänge der Entwickelung
\[
\begin{array}{lll}
1. & -q^{-\frac{2\partial}{a}} e^{-\frac{2\pi i c}{a}}, & \partial > a,\\
2. & q^{-2}, & \partial < a,\\
3. & q^{-2}\left(1-e^{-\frac{2\pi i c}{a}}\right), & \partial = a.
\end{array}
\]

Nehmen wir zunächst an, $n$ sei kein Quadrat, so kommt der Fall 3. nicht
vor, und es ist
\[
N = \sum_{\partial > a}^{\partial} \frac{\partial}{e}\varphi(e) + \sum_{\partial < a}^{a} \frac{a}{e}\varphi(e),
\]
oder, was dasselbe ist,
\weq{13}{N = 2\sum_{\partial > a}^{\partial} \frac{\partial}{e}\varphi(e).}

$C$ ist nach 1. jedenfalls eine $n$te Einheitswurzel, und da es zugleich
eine rationale Zahl ist, so muß es $=\pm 1$ sein.

Ist sodann $n$ ein Quadrat, so kommen $\varphi(\sqrt{n})$ Faktoren von der
Form 3. vor und es ergibt sich:
\weq{14}{N = 2\sum_{\partial > a}^{\partial} \frac{\partial}{e}\varphi(e) + \varphi(\sqrt{n}).}

$C$ ist hier zwar auch von Null verschieden, aber nicht gleich $\pm 1$.
Den Wert von $C$ brauchen wir in diesem Falle nicht näher zu bestimmen.
(Er ist, wie aus der Kreisteilungstheorie folgt, immer ein Teiler von
$\sqrt{n}$).

Wenn im besondern $n$ eine Primzahl ist, so ist
\weq{15}{N = 2n.}

Man kann über $x$ und $y$ immer so verfügen, daß unter Einhaltung der
Bedingungen § 114 $n = \frac{1}{4}(y^2+mx^2)$ kein Quadrat wird. Zu dem
Ende nehme man $x$ durch 4, $y$ durch 2, aber nicht durch 4 teilbar, und
überdies $y$ ohne ungeraden gemeinsamen Teiler mit $mx^2$ an. Ist dann
$\frac{1}{4}(y^2+mx^2)$ ein Quadrat $M^2$, so
%CONT>>
